\section{2026 Mathematical Physics}

\textit{(6 problems)}

\subsection*{Problem 1.}

\begin{exercise}
\label{physics:2026:1}
Consider a sliding pendulum, consisting of a mass $M$ which can move without friction along a horizontal bar, and which is connected by a massless rod of length $l$ to another mass $m$.

\begin{enumerate}
    \item Derive the angular frequency of small oscillations in the $x-z$ plane.
    \item Next consider the sliding conical (also called spherical) pendulum, for which the mass $M$ can move without friction in the horizontal $x-y$ plane. Consider circular motion of both masses for a fixed, not necessarily small, angle $\theta$. What is the angular frequency as a function of $\theta$? Check your result for small $\theta$.
    \item Now consider the inverted sliding pendulum (which can only move in $x-z$ plane, i.e. back to the set-up of 1). As it starts falling from rest at $\theta=\pi$, there comes a point where the tension $T$ in the rod becomes zero. Find the value of $\theta$ when this happens, assuming for simplicity that $M$ is infinitely large.
    \item Same set-up as 3. What is the tension at $\theta=\frac{\pi}{2}$ and $\theta=\pi$ when the inverted sliding pendulum is falling from rest at $\theta=\pi$ but this time $M$ and $m$ are arbitrary?

    \textit{Hint: For an ellipse given by $\frac{x^2}{a^2}+\frac{y^2}{b^2}=1$, the curvature at the extremum with $x=a$ has a radius $\rho=\frac{b^2}{a}$, and at $y=b$ one has $\rho=\frac{a^2}{b}$.}
\end{enumerate}
\end{exercise}

\begin{solution}
Let $X$ be the horizontal coordinate of the mass $M$, and measure $\theta$ from the downward vertical. The coordinates of the mass $m$ are
\[
    x_m=X+l\sin\theta,\qquad z_m=-l\cos\theta .
\]
The Lagrangian for motion in the $x-z$ plane is
\[
    L=\frac{1}{2}(M+m)\dot X^2+ml\cos\theta\,\dot X\dot\theta
      +\frac{1}{2}ml^2\dot\theta^2+mgl\cos\theta .
\]
The coordinate $X$ is cyclic, so
\[
    P_X=(M+m)\dot X+ml\cos\theta\,\dot\theta
\]
is conserved. For the normal small-oscillation problem we may take the horizontal center of mass to be at rest, so $P_X=0$ and
\[
    \dot X=-\frac{ml\cos\theta}{M+m}\dot\theta .
\]
Substituting this into the kinetic energy gives the effective one-coordinate kinetic energy
\[
    T_{\mathrm{eff}}
    =\frac{1}{2}ml^2\left(1-\frac{m}{M+m}\cos^2\theta\right)\dot\theta^2
    =\frac{1}{2}\frac{ml^2(M+m\sin^2\theta)}{M+m}\dot\theta^2 .
\]
For small oscillations about $\theta=0$,
\[
    T_{\mathrm{eff}}=\frac{1}{2}\frac{ml^2M}{M+m}\dot\theta^2+O(\theta^2\dot\theta^2),
    \qquad
    V=-mgl\cos\theta=-mgl+\frac{1}{2}mgl\,\theta^2+O(\theta^4).
\]
Therefore the small-oscillation angular frequency is
\[
    \omega_0^2=\frac{mgl}{ml^2M/(M+m)}
    =\frac{(M+m)g}{Ml},
    \qquad
    \omega_0=\sqrt{\frac{(M+m)g}{Ml}} .
\]

For the conical pendulum, let the horizontal relative vector from $M$ to $m$ have length $l\sin\theta$ and rotate with angular velocity $\Omega$. After eliminating the free horizontal center-of-mass motion, the part of the kinetic energy depending on this rotation is
\[
    T_{\phi}=\frac{1}{2}\mu l^2\sin^2\theta\,\Omega^2,
    \qquad
    \mu=\frac{Mm}{M+m}.
\]
At fixed $\theta$, the effective Lagrangian is
\[
    L_{\mathrm{circ}}=\frac{1}{2}\mu l^2\sin^2\theta\,\Omega^2+mgl\cos\theta .
\]
The condition that $\theta$ remain constant is
\[
    \frac{\partial L_{\mathrm{circ}}}{\partial\theta}
    =\mu l^2\Omega^2\sin\theta\cos\theta-mgl\sin\theta=0 .
\]
For $\sin\theta\neq 0$ this gives
\[
    \Omega^2=\frac{mg}{\mu l\cos\theta}
    =\frac{(M+m)g}{Ml\cos\theta}.
\]
When $\theta$ is small, $\cos\theta\simeq 1$, and hence
\[
    \Omega\simeq \sqrt{\frac{(M+m)g}{Ml}},
\]
which agrees with the small-oscillation result in part 1.

For part 3, take $M=\infty$, so the pivot is fixed. The energy of the mass $m$ released from rest at $\theta=\pi$ is
\[
    E=mgl.
\]
At a later angle,
\[
    \frac{1}{2}ml^2\dot\theta^2-mgl\cos\theta=mgl,
\]
so
\[
    l\dot\theta^2=2g(1+\cos\theta).
\]
The radial force equation is
\[
    T=m(l\dot\theta^2+g\cos\theta).
\]
The tension vanishes when
\[
    2g(1+\cos\theta)+g\cos\theta=0,
\]
or
\[
    \cos\theta=-\frac{2}{3}.
\]
Thus
\[
    \theta=\arccos\left(-\frac{2}{3}\right).
\]

For part 4, keep $M$ finite and again set the conserved horizontal momentum to zero. The effective kinetic energy derived above gives conservation of energy:
\[
    \frac{1}{2}ml^2\left(1-\frac{m}{M+m}\cos^2\theta\right)\dot\theta^2
    -mgl\cos\theta=mgl .
\]
At $\theta=\pi/2$ this becomes
\[
    \frac{1}{2}ml^2\dot\theta^2=mgl,
    \qquad
    \dot\theta^2=\frac{2g}{l}.
\]
The equation of motion for $M$ is
\[
    M\ddot X=T\sin\theta .
\]
Since
\[
    \dot X=-\frac{ml\cos\theta}{M+m}\dot\theta,
\]
we have, at $\theta=\pi/2$,
\[
    \ddot X=\frac{ml}{M+m}\dot\theta^2.
\]
Therefore
\[
    T\left(\frac{\pi}{2}\right)
    =M\ddot X
    =\frac{Mm}{M+m}l\dot\theta^2
    =\frac{2Mm}{M+m}g .
\]
At $\theta=\pi$ the system is instantaneously at rest in the limiting initial configuration. The rod is vertical and the pivot has no horizontal acceleration, so the vertical force balance on $m$ gives
\[
    -T-mg=0,
    \qquad
    T(\pi)=-mg .
\]
The negative sign means that the rod is in compression at the inverted point; a string would not be able to sustain this force.
\end{solution}

\subsection*{Problem 2.}

\begin{exercise}
\label{physics:2026:2}
Two identical circular coils of wire of radius $a$ are separated by a length $h$. The coils each carry a slowly varying sinusoidal current $I(t)=I_0\cos(\omega t)$. The axis of the coils is aligned with the $z$-axis and the geometry is centered at $z=0$.

\begin{enumerate}
    \item At lowest order in the frequency, a magnetostatic approximation is valid. Using this approximation, show that close to the axis, and near $z=0$, the Taylor series for the axial and radial components of the slightly off-axis magnetic field take the approximate form:
    \begin{equation}
        B_z\simeq B_0+\frac{1}{2}\beta_2\left(z^2-\frac{\rho^2}{2}\right)+\cdots, \tag{1}
    \end{equation}
    \begin{equation}
        B_\rho\simeq -\frac{1}{2}\beta_2z\rho+\cdots, \tag{2}
    \end{equation}
    where $B_0$ and $\beta_2$ are determined by the Taylor series of the magnetic field on the $z$-axis:
    \begin{equation}
        B_z(z)\simeq B_0+\frac{1}{2}\beta_2z^2+\cdots. \tag{3}
    \end{equation}
    Here $\rho=\sqrt{x^2+y^2}$.
    \item Using the magnetostatic approximation, determine the magnetic field in the $z$ direction close to the axis of the solenoid, and near $z=0$, to quadratic order in $z$ and $\rho$. Describe the magnetic field when $h=a$.
    \item Determine the electric field close to the axis of the solenoid at $z=0$ to the lowest non-trivial order in the frequency and $\rho$.
    \item Briefly answer the following:
    \begin{enumerate}
        \item For parts 2 and 3, give an estimate for the size of finite-frequency corrections.
        \item For part 2, estimate at what large $z$ the magnetostatic approximation breaks down when computing the magnetic field on the $z$-axis.
    \end{enumerate}
\end{enumerate}
\end{exercise}

\begin{solution}
I use SI units in the formulas below. In a system of units with a different convention for currents, the same derivation applies after replacing the overall factor $\mu_0 I$ by the corresponding magnetostatic normalization.

Near the axis and away from the wires there is no current, so the magnetostatic field obeys
\[
    \nabla\cdot\mathbf B=0,\qquad \nabla\times\mathbf B=0 .
\]
By axial symmetry, write the lowest terms as
\[
    B_z=B_0+\frac{1}{2}\beta_2 z^2+A\rho^2+\cdots,
    \qquad
    B_\rho=Cz\rho+\cdots .
\]
The $\phi$ component of $\nabla\times\mathbf B=0$ gives
\[
    \partial_zB_\rho-\partial_\rho B_z=0,
    \qquad
    C\rho-2A\rho=0,
\]
so $C=2A$. The divergence equation gives
\[
    \frac{1}{\rho}\partial_\rho(\rho B_\rho)+\partial_zB_z
    =2Cz+\beta_2z=0,
\]
hence $C=-\beta_2/2$ and $A=-\beta_2/4$. Therefore
\[
    B_z\simeq B_0+\frac{1}{2}\beta_2\left(z^2-\frac{\rho^2}{2}\right),
    \qquad
    B_\rho\simeq -\frac{1}{2}\beta_2z\rho,
\]
as required.

The on-axis magnetic field of one circular coil centered at $z=z_0$ is
\[
    B_z^{\mathrm{loop}}(z)
    =\frac{\mu_0 I(t)a^2}{2\left(a^2+(z-z_0)^2\right)^{3/2}} .
\]
For two identical coils at $z=\pm h/2$,
\[
    B_z(z,t)
    =\frac{\mu_0 I(t)a^2}{2}
    \left[
      \frac{1}{\left(a^2+(z-h/2)^2\right)^{3/2}}
      +\frac{1}{\left(a^2+(z+h/2)^2\right)^{3/2}}
    \right].
\]
The field is even in $z$, so the linear term vanishes. At $z=0$,
\[
    B_0(t)=\frac{\mu_0 I(t)a^2}{\left(a^2+h^2/4\right)^{3/2}}.
\]
For $f(u)=(a^2+u^2)^{-3/2}$,
\[
    f''(u)=\frac{3(4u^2-a^2)}{(a^2+u^2)^{7/2}},
\]
and hence
\[
    \beta_2(t)=B_z''(0,t)
    =\frac{3\mu_0 I(t)a^2(h^2-a^2)}{\left(a^2+h^2/4\right)^{7/2}} .
\]
Thus, to quadratic order near the center,
\[
    B_z(\rho,z,t)
    \simeq
    \frac{\mu_0 I(t)a^2}{\left(a^2+h^2/4\right)^{3/2}}
    +\frac{3\mu_0 I(t)a^2(h^2-a^2)}{2\left(a^2+h^2/4\right)^{7/2}}
       \left(z^2-\frac{\rho^2}{2}\right).
\]
When $h=a$, the quadratic coefficient $\beta_2$ vanishes. This is the Helmholtz-coil condition: the field near the center is especially uniform, with
\[
    B_z\simeq \frac{8\mu_0 I(t)}{5\sqrt{5}\,a}
\]
up to higher-order spatial corrections.

At $z=0$, Faraday's law gives an azimuthal electric field. For a circle of radius $\rho$ centered on the axis,
\[
    2\pi\rho\,E_\phi(\rho,t)
    =-\frac{d}{dt}\int_0^\rho B_z(\rho',0,t)\,2\pi\rho'\,d\rho' .
\]
The leading non-trivial term in $\rho$ keeps only $B_0(t)$ in the flux:
\[
    E_\phi(\rho,t)
    =-\frac{\rho}{2}\dot B_0(t)+O(\rho^3\dot\beta_2).
\]
Since $I(t)=I_0\cos(\omega t)$,
\[
    E_\phi(\rho,t)
    =
    \frac{\mu_0 I_0\omega a^2}{2\left(a^2+h^2/4\right)^{3/2}}\,
    \rho\sin(\omega t)
    +O(\omega\rho^3\beta_2).
\]
The direction is the azimuthal direction fixed by Lenz's law.

The finite-frequency corrections are controlled by retardation. If $L$ denotes a typical size such as $a$, $h$, $\rho$, or $|z|$, the quasistatic magnetic field in part 2 has relative corrections of order
\[
    O\left(\frac{\omega^2L^2}{c^2}\right),
\]
while the induced electric field in part 3 is already first order in $\omega$ and receives relative corrections of the same retardation size. On the axis at very large $z$, the static dipole field decays as $z^{-3}$, whereas the radiative field decays as $z^{-1}$. The magnetostatic approximation therefore breaks down when
\[
    z\sim \frac{c}{\omega},
\]
up to constants of order one.
\end{solution}

\subsection*{Problem 3.}

\begin{exercise}
\label{physics:2026:3}
Consider a one-dimensional non-relativistic particle of mass $m$ and kinetic energy $E$ scattering off the potential $U(x)$ composed of two static delta-functions
\begin{equation}
    U(x)=\beta[\delta(x)+\delta(x-a)]. \tag{4}
\end{equation}

\begin{enumerate}
    \item Can the particle tunnel through this barrier without reflection? Explain your answer.
    \item If so, at what value(s) of the kinetic energy does this happen?
\end{enumerate}
Now consider a three-dimensional non-relativistic particle of mass $m$ and kinetic energy $E$ scattering off the potential $U(\vec r)$ composed of two static delta-functions
\begin{equation}
    U(\vec r)=\beta_3\left(\delta^{(3)}(\vec r-\vec r_1)+\delta^{(3)}(\vec r-\vec r_2)\right). \tag{5}
\end{equation}
Suppose that for each of the delta-functions in the above equation the S-wave scattering length $a<0$.
\begin{enumerate}
    \setcounter{enumi}{2}
    \item Write explicitly the S-wave bound wave-function near each of the centers. Check that each does not support a bound state for $a<0$.
    \item Can the potential $U(\vec r)$ with two centers support a bound state? If so, under what conditions? Explain your answer.
\end{enumerate}
\textit{Hint: For a hard core potential of radius $R$, the S-wave scattering length $a$ is defined as $d\log\chi(R)/dR=-1/a$, with $\chi(r)$ the S-wave reduced wave-function $\chi(r)\equiv r\psi(r)$. If $\psi(r)$ corresponds to a bound state, think about the constraints on the scattering length $a$.}
\end{exercise}

\begin{solution}
Let
\[
    k=\frac{\sqrt{2mE}}{\hbar},
    \qquad
    \lambda=\frac{2m\beta}{\hbar^2}.
\]
Across a delta-function of strength $\beta$, the wave function is continuous and its derivative jumps by
\[
    \psi'(x_0^+)-\psi'(x_0^-)=\lambda\psi(x_0).
\]
To test whether reflectionless transmission is possible, impose no reflected wave to the left:
\[
    \psi(x)=e^{ikx}\quad (x<0),
\]
and write
\[
    \psi(x)=Ae^{ikx}+Be^{-ikx}\quad (0<x<a),
    \qquad
    \psi(x)=te^{ikx}\quad (x>a).
\]
The matching conditions at $x=0$ give
\[
    A+B=1,
    \qquad
    ik(A-B)-ik=\lambda,
\]
so
\[
    A=1-\frac{i\lambda}{2k},
    \qquad
    B=\frac{i\lambda}{2k}.
\]
The matching conditions at $x=a$ are consistent with zero reflection precisely when
\[
    e^{2ika}
    =-\frac{1+i\lambda/(2k)}{1-i\lambda/(2k)} .
\]
The right hand side has unit modulus, so there are real solutions for $k$. Equivalently,
\[
    ka=\frac{\pi}{2}+\arctan\left(\frac{\lambda}{2k}\right)+n\pi,
    \qquad n=0,1,2,\ldots .
\]
In terms of $\beta$ this can be written as
\[
    \tan(ka)=-\frac{2k}{\lambda}
    =-\frac{\hbar^2k}{m\beta}.
\]
Thus resonant tunneling without reflection is possible at the discrete kinetic energies
\[
    E_n=\frac{\hbar^2k_n^2}{2m},
\]
where $k_n>0$ solves the preceding transcendental equation. Physically, the region between the two delta barriers acts as a Fabry-Perot cavity; at resonance the multiple reflected amplitudes cancel on the incident side.

For the three-dimensional contact problem, write the separation between the two centers as
\[
    d=|\vec r_1-\vec r_2|.
\]
For a single center, an S-wave bound-state solution outside the short-range core has the form
\[
    \psi(r)=C\frac{e^{-\kappa r}}{r},
    \qquad
    E=-\frac{\hbar^2\kappa^2}{2m},
    \qquad
    \kappa>0.
\]
Near the center,
\[
    \psi(r)=C\left(\frac{1}{r}-\kappa+O(r)\right).
\]
The Bethe-Peierls boundary condition encoded by the scattering length is
\[
    \psi(r)\sim C\left(\frac{1}{r}-\frac{1}{a}\right).
\]
For one center this would require $\kappa=1/a$. If $a<0$, then $1/a<0$, which cannot equal the positive number $\kappa$. Therefore one isolated center with negative S-wave scattering length does not support an S-wave bound state.

For two centers, a bound-state ansatz outside the two contact cores is
\[
    \psi(\vec r)
    =
    A_1\frac{e^{-\kappa|\vec r-\vec r_1|}}{|\vec r-\vec r_1|}
    +
    A_2\frac{e^{-\kappa|\vec r-\vec r_2|}}{|\vec r-\vec r_2|}.
\]
Near $\vec r_1$ this becomes
\[
    \psi
    =
    A_1\left(\frac{1}{r_1}-\kappa+O(r_1)\right)
    +A_2\frac{e^{-\kappa d}}{d}+O(r_1),
\]
and similarly near $\vec r_2$. Matching the constant term to $-A_i/a$ gives
\[
    \left(\kappa-\frac{1}{a}\right)A_1
    =\frac{e^{-\kappa d}}{d}A_2,
    \qquad
    \left(\kappa-\frac{1}{a}\right)A_2
    =\frac{e^{-\kappa d}}{d}A_1 .
\]
Nonzero coefficients require
\[
    \left(\kappa-\frac{1}{a}\right)^2
    =\frac{e^{-2\kappa d}}{d^2}.
\]
The antisymmetric solution would require
\[
    \kappa-\frac{1}{a}=-\frac{e^{-\kappa d}}{d},
\]
which is impossible for $a<0$ because the left hand side is positive. The symmetric solution requires
\[
    \kappa-\frac{1}{a}=\frac{e^{-\kappa d}}{d}.
\]
For $a<0$, this is
\[
    \kappa+\frac{1}{|a|}=\frac{e^{-\kappa d}}{d}.
\]
The left hand side increases with $\kappa$, while the right hand side decreases with $\kappa$. A positive solution exists exactly when the equation is satisfied at threshold in the correct direction:
\[
    \frac{1}{|a|}<\frac{1}{d},
    \qquad\text{or}\qquad
    d<|a|.
\]
Thus two individually unbound centers can support one symmetric bound state if they are closer than the magnitude of the negative scattering length. At $d=|a|$ the bound state is at threshold, and for $d>|a|$ there is no bound state.
\end{solution}

\subsection*{Problem 4.}

\begin{exercise}
\label{physics:2026:4}
\begin{enumerate}
    \item Write down the condition for thermodynamic equilibrium between the liquid and gas phases along a liquid-gas coexistence curve.
    \item Using your solution to part 1 and taking into account the entropy change at a liquid-gas phase transition, derive the relation for the vapor pressure along a liquid-gas coexistence curve.
\end{enumerate}
\end{exercise}

\begin{solution}
Along a liquid-gas coexistence curve the two phases can exchange particles, volume, and energy without any net thermodynamic drive. Therefore their temperature, pressure, and chemical potential agree:
\[
    T_l=T_g=T,\qquad P_l=P_g=P,\qquad \mu_l(T,P)=\mu_g(T,P).
\]
Equivalently, the molar Gibbs free energies of the two phases are equal.

Differentiate the coexistence condition:
\[
    d\mu_l=d\mu_g.
\]
For one particle, or per mole with the corresponding molar quantities, the Gibbs-Duhem relation gives
\[
    d\mu=-s\,dT+v\,dP,
\]
where $s$ is entropy per particle and $v$ is volume per particle. Hence
\[
    -s_l\,dT+v_l\,dP=-s_g\,dT+v_g\,dP.
\]
Rearranging,
\[
    \frac{dP}{dT}
    =\frac{s_g-s_l}{v_g-v_l}.
\]
At a first-order liquid-gas transition, the latent heat per particle is
\[
    L=T(s_g-s_l).
\]
Therefore the exact Clapeyron equation is
\[
    \frac{dP}{dT}
    =\frac{L}{T(v_g-v_l)}.
\]
For vapor pressure away from the critical point, one usually has $v_g\gg v_l$ and the gas is approximately ideal:
\[
    v_g\simeq \frac{k_BT}{P}.
\]
Then
\[
    \frac{dP}{dT}\simeq \frac{LP}{k_BT^2},
    \qquad
    \frac{d\log P}{dT}\simeq \frac{L}{k_BT^2}.
\]
If $L$ is approximately constant over the temperature range of interest, integration gives the Clausius-Clapeyron vapor-pressure relation
\[
    \log P(T)=C-\frac{L}{k_BT},
\]
or, relative to a reference point $(T_0,P_0)$,
\[
    P(T)=P_0
    \exp\left[
        -\frac{L}{k_B}\left(\frac{1}{T}-\frac{1}{T_0}\right)
    \right].
\]
For molar quantities, replace $k_B$ by the gas constant $R$ and let $L$ be the molar latent heat.
\end{solution}

\subsection*{Problem 5.}

\begin{exercise}
\label{physics:2026:5}
Consider the $\phi^4$ model with a real scalar field $\phi(x)$ in $3+1$ dimensional Minkowski spacetime with metric $(-,+,+,+)$. Its Lagrangian density is
\begin{equation}
    \mathcal L=-\frac{1}{2}\partial_\mu\phi\partial^\mu\phi-\frac{1}{2}m^2\phi^2-\frac{1}{4!}g\phi^4 . \tag{6}
\end{equation}

\begin{enumerate}
    \item Write down the propagator and the interaction vertex for this model in momentum space.
    \item Compute the one-loop correction to the $4$-point function using dimensional regularization.
    \item Determine the counterterm needed to cancel the divergence and write the renormalized coupling at $1$-loop order (You may use minimal subtraction).
    \item Write down the renormalization group flow equation for coupling constant.
\end{enumerate}
You may find the following formula useful
\begin{equation}
    (AB)^{-1}=\int_0^1dx[xA+(1-x)B]^{-2}, \tag{7}
\end{equation}
\begin{equation}
    \int d^dk\frac{1}{(-k^2-2p\cdot k-M^2+i\epsilon)^s}
    =
    (-1)^s i\pi^{d/2}
    \frac{\Gamma(s-d/2)}{\Gamma(s)}
    (-p^2+M^2-i\epsilon)^{d/2-s}, \tag{8}
\end{equation}
where $\Gamma(z)$ is the Gamma function which has a simple pole at the origin.
\begin{equation}
    \Gamma(z)=\frac{1}{z}-\gamma+O(z), \tag{9}
\end{equation}
where $\gamma$ is the Euler constant.
\end{exercise}

\begin{solution}
With metric $(-,+,+,+)$, the free propagator is
\[
    \frac{i}{-p^2-m^2+i\epsilon}.
\]
The interaction term is $-g\phi^4/4!$, so the four-point vertex is
\[
    -ig.
\]

Let the external momenta be incoming and define the Mandelstam-channel momenta
\[
    P_s=p_1+p_2,\qquad P_t=p_1+p_3,\qquad P_u=p_1+p_4 .
\]
The one-loop correction is the sum of the three bubble diagrams. For one channel with momentum $P$, the contribution is
\[
    \mathcal M_1(P)
    =
    \frac{(-ig)^2}{2}
    \int\frac{d^dk}{(2\pi)^d}
    \frac{i}{-k^2-m^2+i\epsilon}
    \frac{i}{-(k+P)^2-m^2+i\epsilon}.
\]
The factor $1/2$ is the symmetry factor of the bubble. Using a Feynman parameter,
\[
    \frac{1}{AB}=\int_0^1dx\,[xA+(1-x)B]^{-2},
\]
and shifting the loop momentum gives
\[
    \mathcal M_1(P)
    =
    \frac{g^2}{2}
    \int_0^1dx
    \int\frac{d^dq}{(2\pi)^d}
    \frac{1}{\left[-q^2-\Delta(P,x)+i\epsilon\right]^2},
\]
where
\[
    \Delta(P,x)=m^2+x(1-x)P^2.
\]
Applying the dimensional-regularization formula in $d=4-\epsilon$ gives
\[
    \mathcal M_1(P)
    =
    \frac{ig^2}{32\pi^2}
    \left[
        \frac{2}{\epsilon}
        -\gamma+\log(4\pi)
        -\int_0^1dx\,
        \log\frac{\Delta(P,x)-i\epsilon}{\mu^2}
    \right]
    +O(\epsilon),
\]
where the usual scale $\mu$ has been introduced to keep $g$ dimensionless. The divergent part of one channel is therefore
\[
    \mathcal M_1(P)\big|_{\mathrm{div}}
    =
    \frac{ig^2}{16\pi^2}\frac{1}{\epsilon}.
\]
There are three channels, so
\[
    \mathcal M_{\mathrm{1-loop}}\big|_{\mathrm{div}}
    =
    \frac{3ig^2}{16\pi^2}\frac{1}{\epsilon}.
\]

Introduce a counterterm
\[
    \mathcal L_{\mathrm{ct}}=-\frac{\delta g}{4!}\phi^4,
\]
whose vertex is $-i\delta g$. Minimal subtraction cancels the pole by choosing
\[
    \delta g=\frac{3g^2}{16\pi^2}\frac{1}{\epsilon}.
\]
Equivalently, the bare coupling is
\[
    g_0=\mu^\epsilon\left(g+\frac{3g^2}{16\pi^2}\frac{1}{\epsilon}+O(g^3)\right).
\]
After subtraction, the renormalized four-point function can be written as
\[
    \Gamma_{\mathrm{ren}}^{(4)}
    =
    -ig
    +
    \frac{ig^2}{32\pi^2}
    \sum_{P=P_s,P_t,P_u}
    \left[
        -\int_0^1dx\,
        \log\frac{m^2+x(1-x)P^2-i\epsilon}{\mu^2}
    \right]
    +O(g^3),
\]
up to the finite convention-dependent constants removed together with the pole if one uses $\overline{\mathrm{MS}}$ instead of MS.

Because the bare coupling is independent of the renormalization scale, differentiating $g_0$ with respect to $\mu$ gives the one-loop beta function
\[
    \mu\frac{dg}{d\mu}
    =
    \beta(g)
    =
    \frac{3g^2}{16\pi^2}+O(g^3).
\]
The positive sign means that the coupling grows toward the ultraviolet at this order.
\end{solution}

\subsection*{Problem 6.}

\begin{exercise}
\label{physics:2026:6}
The Klein disk model is an analytical description of non-Euclidean geometry in terms of coordinates $(x^1,x^2)$ for points inside a unit disk. The distance $d(\vec x,\vec y)$ between a point with coordinates $\vec x=(x^1,x^2)$ and another point with coordinates $\vec y=(y^1,y^2)$, both inside the unit disk, is defined to be
\begin{equation}
    \cosh d(\vec x,\vec y)
    =
    \frac{1-\vec x\cdot\vec y}{\sqrt{1-\vec x\cdot\vec x}\sqrt{1-\vec y\cdot\vec y}}, \tag{10}
\end{equation}
where $\vec x\cdot\vec y=x^1y^1+x^2y^2$.

\begin{enumerate}
    \item Compute the metric with components $g_{11}$, $g_{12}$ and $g_{22}$ from the formula for $d$ by considering $\vec y$ near $\vec x$ (so $y^j=x^j+dx^j$, $j=1,2$).
    \item Introduce the standard polar coordinates $(r,\theta)$ and write the line elements $ds^2=g_{\mu\nu}dx^\mu dx^\nu$ in polar coordinates.
    \item Determine the nonvanishing Christoffel symbols in polar coordinates.
    \item Compute the component $R_{r\theta r\theta}$ of the Riemann tensor.
    \item Show that for the Klein disk model, the Riemann curvature has the form
    \begin{equation}
        R_{\mu\nu\rho\sigma}
        =
        K(g_{\mu\sigma}g_{\nu\rho}-g_{\mu\rho}g_{\nu\sigma}), \tag{11}
    \end{equation}
    and determine the value of $K$.
    \item Compute the scalar curvature.
\end{enumerate}
\end{exercise}

\begin{solution}
Set
\[
    r^2=\vec x\cdot\vec x,
    \qquad
    \vec y=\vec x+d\vec x.
\]
Expanding the distance formula to second order and using
\[
    \cosh ds=1+\frac{1}{2}ds^2+O(ds^4),
\]
one obtains
\[
    ds^2
    =
    \frac{(1-r^2)\,d\vec x\cdot d\vec x+(\vec x\cdot d\vec x)^2}{(1-r^2)^2}.
\]
Therefore
\[
    g_{ij}
    =
    \frac{(1-r^2)\delta_{ij}+x_ix_j}{(1-r^2)^2}.
\]
In components,
\[
    g_{11}
    =
    \frac{1-(x^2)^2}{(1-r^2)^2},
    \qquad
    g_{22}
    =
    \frac{1-(x^1)^2}{(1-r^2)^2},
    \qquad
    g_{12}=g_{21}
    =
    \frac{x^1x^2}{(1-r^2)^2}.
\]

In polar coordinates,
\[
    d\vec x\cdot d\vec x=dr^2+r^2d\theta^2,
    \qquad
    \vec x\cdot d\vec x=r\,dr.
\]
Thus
\[
    ds^2
    =
    \frac{dr^2}{(1-r^2)^2}
    +\frac{r^2}{1-r^2}d\theta^2.
\]
Hence
\[
    g_{rr}=\frac{1}{(1-r^2)^2},
    \qquad
    g_{\theta\theta}=\frac{r^2}{1-r^2},
    \qquad
    g_{r\theta}=0.
\]

For a diagonal metric of the form
\[
    ds^2=A(r)\,dr^2+B(r)\,d\theta^2,
\]
with
\[
    A(r)=\frac{1}{(1-r^2)^2},
    \qquad
    B(r)=\frac{r^2}{1-r^2},
\]
the nonvanishing Christoffel symbols are
\[
    \Gamma^r_{rr}=\frac{A'}{2A},
    \qquad
    \Gamma^r_{\theta\theta}=-\frac{B'}{2A},
    \qquad
    \Gamma^\theta_{r\theta}=\Gamma^\theta_{\theta r}=\frac{B'}{2B}.
\]
Since
\[
    A'=\frac{4r}{(1-r^2)^3},
    \qquad
    B'=\frac{2r}{(1-r^2)^2},
\]
we get
\[
    \Gamma^r_{rr}=\frac{2r}{1-r^2},
    \qquad
    \Gamma^r_{\theta\theta}=-r,
    \qquad
    \Gamma^\theta_{r\theta}=\Gamma^\theta_{\theta r}
    =\frac{1}{r(1-r^2)}.
\]

Using the convention
\[
    R^\lambda{}_{\mu\rho\nu}
    =
    \partial_\rho\Gamma^\lambda_{\nu\mu}
    -\partial_\nu\Gamma^\lambda_{\rho\mu}
    +\Gamma^\lambda_{\rho\alpha}\Gamma^\alpha_{\nu\mu}
    -\Gamma^\lambda_{\nu\alpha}\Gamma^\alpha_{\rho\mu},
\]
we find
\[
    R^r{}_{\theta r\theta}
    =
    \partial_r\Gamma^r_{\theta\theta}
    +\Gamma^r_{rr}\Gamma^r_{\theta\theta}
    -\Gamma^r_{\theta\theta}\Gamma^\theta_{r\theta}.
\]
Substituting the Christoffel symbols,
\[
    R^r{}_{\theta r\theta}
    =
    -1-\frac{2r^2}{1-r^2}+\frac{1}{1-r^2}
    =
    -\frac{r^2}{1-r^2}.
\]
Lowering the first index gives
\[
    R_{r\theta r\theta}
    =
    g_{rr}R^r{}_{\theta r\theta}
    =
    -\frac{r^2}{(1-r^2)^3}.
\]

Since
\[
    g_{rr}g_{\theta\theta}
    =
    \frac{r^2}{(1-r^2)^3},
\]
the computed component satisfies
\[
    R_{r\theta r\theta}
    =
    -g_{rr}g_{\theta\theta}.
\]
Comparing this with the convention used in the problem,
\[
    R_{\mu\nu\rho\sigma}
    =
    K(g_{\mu\sigma}g_{\nu\rho}-g_{\mu\rho}g_{\nu\sigma}),
\]
and taking $(\mu,\nu,\rho,\sigma)=(r,\theta,r,\theta)$, we get
\[
    R_{r\theta r\theta}
    =
    K(0-g_{rr}g_{\theta\theta}).
\]
Therefore
\[
    K=1
\]
with the sign convention of equation (11). Equivalently, in the more common convention
$R_{\mu\nu\rho\sigma}=K_{\mathrm{std}}(g_{\mu\rho}g_{\nu\sigma}-g_{\mu\sigma}g_{\nu\rho})$, this is constant curvature $K_{\mathrm{std}}=-1$.

Contracting equation (11) in two dimensions gives
\[
    R_{\nu\sigma}=-(2-1)K\,g_{\nu\sigma}=-K g_{\nu\sigma}.
\]
Thus the scalar curvature is
\[
    R=g^{\nu\sigma}R_{\nu\sigma}=-2K=-2.
\]
\end{solution}
